\section*{Введение}
\addcontentsline{toc}{section}{Введение}

Настоящая работа выполняется в рамках группового проекта по созданию платформы персонализированной подготовки к ЕГЭ по математике. 
Платформа строится на пользовательской аналитике: из истории решённых задач формируется вектор компетенций ученика по темам кодификатора ФИПИ, который затем управляет генерацией учебного контента. Фокусом данной работы является модуль генерации конспектов с визуальными материалами.

Типичный конспект по математике без иллюстраций неполноценен для ряда тем ЕГЭ. Задачи на графики функций, взаимное расположение геометрических фигур или анализ производной по графику требуют не просто словесного описания, а корректного визуального представления. Использование моделей генерации изображений по текстовому описанию для этой цели неприемлемо: такие модели систематически нарушают точные геометрические отношения и некорректно воспроизводят подписи и формулы ~\cite{Bosheah2025, Wang2025MagicGeo}. В данной работе предлагается двухшаговый процесс: языковая модель определяет структуру конспекта и места для визуальных элементов, формирует их текстовые описания, затем отдельный модуль на основе этих описаний генерирует JSON-спецификацию, а специализированный движок рендеринга строит по ней изображение в формате масштабируемой векторной графики (Scalable Vector Graphics, SVG). Такой подход обеспечивает геометрическую корректность, валидируемость выхода и возможность адаптировать сложность иллюстрации под уровень ученика.

Цель работы --- разработать и реализовать модуль генерации персонализированных учебных конспектов по математике ЕГЭ, включающих корректные визуальные материалы, на основе пользовательской аналитики и языковых моделей.

Для достижения этой цели в работе решаются следующие задачи.

\begin{enumerate}
    \item Проанализировать существующие подходы к персонализации учебного контента, генерации учебных материалов и автоматическому созданию математических визуализаций.
    \item Спроектировать архитектуру процесса генерации конспектов: от вектора компетенций ученика к структурированному тексту с визуальными элементами.
    \item Разработать схему JSON-представления для математических визуализаций и реализовать движок SVG-рендеринга на её основе.
    \item Интегрировать модуль визуализаций в общий процесс генерации конспектов с учётом профиля ученика.
    \item Оценить качество генерируемых конспектов.
\end{enumerate}

Разработанный модуль встраивается в действующую платформу подготовки к ЕГЭ «Прорыв» ~\cite{Proriv} и обеспечивает автоматическое создание персонализированных учебных материалов для реальных пользователей. Предложенный подход к структурированной генерации визуализаций через промежуточное JSON-представление применим и для других предметов, где требуются точные графические иллюстрации.

\paragraph{Разделение труда в команде.} Работа выполнена совместно двумя авторами. А.~С.~Низов разрабатывал двухпроходный планировщик визуалов конспекта, педагогическую таксономию типов иллюстраций, ветку пошаговой геометрии и стереометрии, а также проводил количественные эксперименты на конспектном и геометрическом корпусах. Д.~А.~Ковырков разрабатывал конвертер плана визуалов в описание сцены, движок векторного рендеринга со схемами Pydantic и решателем геометрических ограничений, а также инженерные итерации улучшений прототипа и инфраструктуру экспериментального конвейера. Промежуточный JSON\nobreakdash-формат описания сцены и схема журналирования прогона сопровождались авторами совместно.
